\VerdierPhysicalPageAnchor{0035}%
\VerdierSemanticAnchor{ch1-1-1-1-fiber-category-cont}%
thus the subcategory of $\mathcal C$ with the same set of objects as
$\mathcal C$ and whose morphisms are the morphisms of $\mathcal C$ of
degree $e$.

\VerdierSemanticAnchor{ch1-1-1-1-graded-functor}
Let $(\mathcal C,G)$ and $(\mathcal C',G)$ be two $G$-graded categories.
A \emph{graded functor} from $(\mathcal C,G)$ to $(\mathcal C',G)$ is a
cartesian functor $F$ over $\mathcal G$. We denote by
$F_e\colon\mathcal C_e\to\mathcal C'_e$ the functor obtained by restriction
to the fiber categories.
\VerdierSemanticAnchor{ch1-1-1-1-compatible-morphism}
Finally, a morphism of graded functors $m\colon F\to F'$ is said to be
\emph{compatible with the grading} if it is a
$\mathcal G$-morphism of functors.

\VerdierSemanticAnchor{ch1-1-1-2}
\noindent\textbf{1.1.2.}\quad
Let $\mathcal C$ be a $G$-graded category and let $g$ be an element of
$G$. For every pair $X,Y$ of objects of $\mathcal C$, let
\VerdierSemanticAnchor{ch1-1-1-2-homogeneous-hom}%
$\operatorname{Hom}^{g}(X,Y)$ denote the set of morphisms from $X$ to $Y$
of degree $g$. Then:
\VerdierSemanticAnchor{ch1-1-1-2-eq1}
\begingroup
\makeatletter
\tagsleft@true
\let\veqno\@@leqno
\makeatother
\begin{equation}
  \operatorname{Hom}(X,Y)
  =\coprod_{g\in G}\operatorname{Hom}^{g}(X,Y).
  \tag{1.1.2.1}
\end{equation}
\endgroup

\VerdierSemanticAnchor{ch1-1-1-2-graded-hom-set}
Thus the set $\operatorname{Hom}(X,Y)$ is endowed with the structure of a
$G$-graded set. Since degree is functorial, for every morphism
$m\colon X\to Y$ and every morphism $n\colon Y\to Z$ we have:
\VerdierSemanticAnchor{ch1-1-1-2-eq2}
\begin{equation}
  \deg(n)\deg(m)=\deg(nm).
  \tag{1.1.2.2}
\end{equation}

\VerdierSemanticAnchor{ch1-1-1-2-fibered-criterion}
The preceding formulas merely express the fact that $\mathcal C$ is a
category over $\mathcal G$. Let us now express that the category
$\mathcal C$ is fibered over $\mathcal G$. For this it is enough to say
that every object of $\mathcal C$ is the target of a hypercartesian
morphism of prescribed degree [19]. But one immediately checks that the
hypercartesian morphisms of $\mathcal C$ are the isomorphisms of
$\mathcal C$, and consequently that $\mathcal C$ is fibered if and only if
it has the following property:

\VerdierSemanticAnchor{ch1-1-1-2-eq3}
\noindent\textnormal{(1.1.2.3)}\quad
\begin{minipage}[t]{0.78\textwidth}
\emph{For every element $g$ of $G$ and every object $X$ of $\mathcal C$,
there exists an isomorphism with target $X$ and degree $g$.}
\end{minipage}
\par\medskip

\VerdierSemanticAnchor{ch1-1-1-2-degree-preserving-functors}
Now let $(\mathcal C,G)$ and $(\mathcal C',G)$ be two $G$-graded
categories and let $F\colon\mathcal C\to\mathcal C'$ be a graded functor.
The functor $F$ is first of all a $\mathcal G$-functor and hence preserves
the degrees of morphisms. It also sends cartesian morphisms of $\mathcal C$
to cartesian morphisms of $\mathcal C'$, which is automatic because the
cartesian morphisms of $\mathcal C$ are isomorphisms. Thus the graded
functors from $\mathcal C$ to $\mathcal C'$ are simply
\emph{the functors from $\mathcal C$ to $\mathcal C'$ that preserve
degree}.

\VerdierSemanticAnchor{ch1-1-1-2-degree-e-morphism}
Finally, a morphism of graded functors $m\colon F\to F'$ compatible with
the grading is simply a morphism of functors of degree $e$, \emph{i.e.}
the datum, for every object $X$ of $\mathcal C$, of a morphism:
\VerdierSemanticAnchor{ch1-1-1-2-component}
\[
  m(X)\colon F(X)\longrightarrow F'(X)
\]
